\chapter{The USD asset files}
\label{ch:usd}

\textbf{Role:} the car as Isaac Lab actually loads it. These files are
generated (Chapter~\ref{ch:asset}), so they are explained block by block
rather than line by line; every block is covered. Listings are the committed
files. The paths are relative to \srcpath{assets/}.

\newcommand{\usdpart}[3]{\lstinputlisting[style=plain,basicstyle=\ttfamily\footnotesize,firstline=#2,lastline=#3,firstnumber=#2]{#1}}
\newcommand{\usdfile}[1]{\lstinputlisting[style=plain,basicstyle=\ttfamily\footnotesize]{#1}}

\section{\texttt{tricycle/tricycle.usda}: the entry point}

\usdfile{\srcroot/assets/tricycle/tricycle.usda}

\begin{explain}
\lis{1}{18} Layer metadata. \code{creator} records the converter version;
  \code{doc} keeps the temp paths the converter went through (three
  stages: MuJoCo USD Converter, the \code{usdex} rewrite, the Isaac asset
  transformer); \code{kilogramsPerUnit}/\code{metersPerUnit} say the file is
  in SI units; \code{upAxis = "Z"} matches Isaac Lab and MuJoCo.
  \code{defaultPrim = "tricycle"} tells a referencing stage which prim to
  graft in, which is how \code{UsdFileCfg} places \srcpath{/tricycle} at
  \code{{ENV_REGEX_NS}/Tricycle}.
\lis{20}{27} The root prim, an \code{Xform}, pulls in \srcpath{payloads/base.usda}
  by reference and declares a variant set named \code{Physics} whose
  selected variant is \code{"physics"}.
\lis{28}{48} The four variants. \code{mujoco}, \code{physics} and
  \code{physx} each add a payload file; \code{none} adds nothing (a
  visual-only car). Only the selected variant's payload is composed, so with
  \code{"physics"} selected, \srcpath{Physics/physics.usda} is loaded and the
  MuJoCo/PhysX-specific files are not.
\end{explain}

\section{\texttt{payloads/base.usda}: geometry}

\usdpart{\srcroot/assets/tricycle/payloads/base.usda}{1}{46}

\begin{explain}
\lis{1}{22} Metadata as above, plus \code{subLayers = [@./robot.usda@]}:
  the robot-schema layer is stacked under this one, so its \code{over}
  blocks apply to the prims defined here.
\lis{23}{30} The root prim gets \code{GeomModelAPI} and
  \code{IsaacRobotAPI}; \code{kind = "component"} marks it as one model.
\li{31} \code{extentsHint}: the bounding box of the whole car, from the back
  of the rear wheels ($x=-0.25$) to the front of the front wheel
  ($x=0.32$), $\pm0.26$ in $y$, and $0$ to $0.2$ in $z$ (wheel bottoms to
  wheel tops). Unchanged by the new body because the frame lies inside this
  box.
\lis{32}{34} Empty robot-joint/robot-link relationships and the robot type,
  filled in by Isaac's asset transformer.
\li{36} All geometry lives under a \code{Geometry} scope.
\lis{38}{45} The chassis body: an \code{Xform} with \code{IsaacLinkAPI},
  translated to $z=0.14$ (the MJCF body \code{pos}), identity orientation,
  unit scale. \code{xformOpOrder} fixes the order translate $\to$ orient
  $\to$ scale.
\end{explain}

\usdpart{\srcroot/assets/tricycle/payloads/base.usda}{47}{130}

\begin{explain}
\lis{47}{66} The left rail as a \code{Cube} prim. USD's \code{Cube} is a
  unit cube of edge \code{size = 2} (so it spans $-1..1$), and the box's
  half-sizes are applied as \code{xformOp:scale = (0.1961, 0.02, 0.015)}.
  \code{xformOp:translate} is the rail centre and \code{xformOp:orient} is
  the yaw as a quaternion in $(w, x, y, z)$ order: $(0.98576, 0, 0,
  -0.16814)$ is a rotation of $-19.4^\circ$ about $z$. The prim carries
  \code{PhysicsCollisionAPI} (it is a collider) and
  \code{MaterialBindingAPI} with the binding to the shared physics material.
  \code{displayColor}/\code{displayOpacity} come from the MJCF
  \code{rgba}. \code{extent} is the unit cube's own bound; the scale is
  applied on top.
\lis{68}{87} The right rail, identical except for the mirrored translate and
  the opposite quaternion sign.
\lis{89}{108} The crossbar: no rotation, half-sizes $(0.02, 0.15, 0.015)$
  at $x=-0.15$.
\lis{110}{129} The deck: half-sizes $(0.06, 0.065, 0.01)$ at $x=-0.08$, in
  the lighter colour.
\end{explain}

\usdpart{\srcroot/assets/tricycle/payloads/base.usda}{131}{194}

\begin{explain}
\lis{131}{138} The fork body under the chassis, translated to $(0.22, 0,
  -0.04)$: the MJCF body \code{pos}.
\lis{140}{161} The fork's capsule. USD capsules are along an axis
  (\code{axis = "Z"}) with \code{height} being the cylinder part (0.04) and
  \code{radius} 0.012. The odd \code{orient = (6.1e-17, 1, 0, 0)} is a
  $180^\circ$ flip about $x$: the MJCF \code{fromto} pointed downwards and
  the converter preserved that direction exactly.
\lis{163}{193} The front wheel body nested inside the fork (at the fork
  origin), containing a \code{Sphere} of radius 0.1, coloured dark grey,
  with the same collision and material tags.
\end{explain}

\usdpart{\srcroot/assets/tricycle/payloads/base.usda}{195}{266}

\begin{explain}
\lis{195}{225} The left rear wheel body at $(-0.15, 0.16, -0.04)$ with its
  sphere.
\lis{226}{255} The right rear wheel body at $(-0.15, -0.16, -0.04)$ with its
  sphere.
\lis{259}{266} A \code{Physics} scope holding the one shared
  \code{PhysicsMaterial} that all colliders bind to. Its attributes are
  added by \code{physics.usda}.
\end{explain}

\begin{isaacbox}{what Newton builds from this}
Reading the composed stage, Isaac Lab's Newton backend creates one Newton
body per prim with \code{PhysicsRigidBodyAPI} (chassis, fork, three wheels),
one Newton shape per prim with \code{PhysicsCollisionAPI} (four cubes, one
capsule, three spheres), and one joint per \code{PhysicsRevoluteJoint}
prim. Body mass and inertia are accumulated from the shapes'
\code{physics:mass} and geometry, exactly as MuJoCo did when it compiled the
MJCF.
\end{isaacbox}

\section{\texttt{payloads/robot.usda}: robot schema tags}

\usdfile{\srcroot/assets/tricycle/payloads/robot.usda}

\begin{explain}
\lis{1}{19} Metadata.
\lis{20}{25} \code{over "tricycle"}: adds \code{IsaacRobotAPI} and the two
  relationship attributes that list the robot's joints and links.
\lis{27}{78} A mirror of the body tree in which every body gets
  \code{IsaacLinkAPI} and every joint gets \code{IsaacJointAPI}. These tags
  are how Isaac Lab tools recognise ``this prim is a robot link/joint''
  independent of the physics variant. Note the joint prims (\code{steer},
  \code{front_wheel}, \code{rear_left}, \code{rear_right}) live under the
  \emph{child} body of the joint, which is also where \code{physics.usda}
  defines them.
\end{explain}

\section{\texttt{payloads/Physics/physics.usda}: bodies, joints, masses}

\usdpart{\srcroot/assets/tricycle/payloads/Physics/physics.usda}{1}{65}

\begin{explain}
\lis{1}{18} Metadata.
\lis{20}{22} Overrides on the existing prims.
\lis{24}{28} The chassis becomes a rigid body (\code{PhysicsRigidBodyAPI}),
  the articulation root (\code{PhysicsArticulationRootAPI},
  \code{NewtonArticulationRootAPI}), and
  \code{newton:selfCollisionEnabled = 0}: bodies of the car do not collide
  with each other. This is why the rails can end inside the wheel spheres
  and the fork capsule can sit inside the front wheel without generating
  contacts.
\lis{30}{37} The left rail gets \code{NewtonCollisionAPI} (Newton-specific
  collision settings; \code{newton:contactGap = 0}) and \code{PhysicsMassAPI}
  with \code{physics:mass = 0.8}. \code{physics:density = 1000} is written
  by the converter but mass takes precedence.
\lis{39}{64} The same for the right rail (0.8), crossbar (0.8) and deck
  (2.6). The four masses sum to 5.0~kg.
\end{explain}

\usdpart{\srcroot/assets/tricycle/payloads/Physics/physics.usda}{66}{125}

\begin{explain}
\lis{66}{68} The fork becomes a rigid body.
\lis{70}{76} Its capsule: collision tags and 0.1~kg.
\lis{78}{80} The front wheel body becomes a rigid body.
\lis{82}{89} Its sphere: collision tags and 0.5~kg.
\lis{92}{105} The \code{front_wheel} joint: a \code{PhysicsRevoluteJoint}
  about axis \code{Y}, between \code{body0} (the fork) and \code{body1} (the
  wheel body), located at the origin of both (\code{localPos0/1 = 0},
  identity rotations). It carries \code{PhysicsDriveAPI:angular} and
  \code{PhysicsJointStateAPI:angular}, the schemas through which Isaac Lab
  sets drive targets/gains and reads joint state. No limits: it spins
  freely.
\lis{108}{124} The \code{steer} joint: revolute about \code{Z}, between the
  chassis (\code{body0}) and the fork (\code{body1}), anchored at
  \code{localPos0 = (0.22, 0, -0.04)} in the chassis and at the fork's
  origin. Limits $\pm34.377^\circ$ are the MJCF's $\pm0.6$~rad converted to
  degrees, which is USD's unit for joint limits.
\end{explain}

\usdpart{\srcroot/assets/tricycle/payloads/Physics/physics.usda}{126}{198}

\begin{explain}
\lis{126}{153} The left rear wheel: rigid body, 0.5~kg sphere, and the
  \code{rear_left} revolute joint about \code{Y} between chassis and wheel,
  anchored at $(-0.15, 0.16, -0.04)$ in the chassis frame. No limits.
\lis{155}{183} The right rear wheel and \code{rear_right} joint, mirrored.
\lis{186}{197} The shared physics material gets \code{PhysicsMaterialAPI} and
  \code{NewtonMaterialAPI} with the friction values from the MJCF default
  geom: dynamic 0.8, torsional 0.005, rolling 0.0001.
\end{explain}

\begin{isaacbox}{how the env cfg's actuators meet these joints}
\code{ImplicitActuatorCfg} entries in \code{tricycle_env_cfg.py} are matched
to these joint prims by name (\code{joint_names_expr}). At start-up Isaac Lab
writes the configured stiffness, damping and limits into the joints' drive
attributes (or, with \srcpath{use_newton_actuators}, into native Newton
actuators). The USD carries no gains of its own, which is why the cfg must
supply them.
\end{isaacbox}

\section{\texttt{payloads/Physics/mujoco.usda}: MuJoCo variant}

\usdfile{\srcroot/assets/tricycle/payloads/Physics/mujoco.usda}

\begin{explain}
\lis{1}{21} Metadata; this layer sublayers \code{physics.usda}, so it is
  ``physics plus MuJoCo extras''.
\lis{23}{121} For every collider, \code{MjcCollisionAPI}; for every joint,
  \code{MjcJointAPI} with \code{mjc:damping = 0.02} (the MJCF default joint
  damping), and the generic drive/state schemas deleted because MuJoCo would
  read its own. This variant would be used to round-trip the asset back into
  a MuJoCo-native pipeline; Isaac Lab's Newton backend uses the
  \code{physics} variant instead, so \emph{the 0.02 joint damping from the
  MJCF is not active during training}. The wheels' damping in training comes
  from the actuator configs (and the free front wheel has none beyond
  solver defaults).
\lis{123}{131} \code{MjcMaterialAPI} on the shared material.
\end{explain}

\section{\texttt{payloads/Physics/physx.usda}: PhysX variant}

\usdfile{\srcroot/assets/tricycle/payloads/Physics/physx.usda}

\begin{explain}
\lis{1}{21} Metadata; also sublayers \code{physics.usda}.
\lis{23}{66} Adds \code{PhysxJointAPI} to the four joints. Not loaded with
  Newton.
\end{explain}

\section{\texttt{config.yaml} and \texttt{.asset\_hash}}

\usdfile{\srcroot/assets/config.yaml}

\begin{explain}
\lis{1}{3} The converter's inputs: source MJCF (the temp file
  \code{write_mjcf} produced), the output directory and the output file name
  (with the subfolder the converter chose).
\li{4} \code{force_usd_conversion: true}: reconvert even if the hash
  matches.
\li{5} \code{make_instanceable: true}: mark the asset so that many copies
  share geometry.
\li{6} Which physics variant to select: \code{physics}.
\lis{7}{9} Mesh handling: keep separate meshes, no collision-from-visuals,
  convex-hull collision for meshes (irrelevant for primitives).
\li{10} \code{self_collision: false}: the source of
  \code{newton:selfCollisionEnabled = 0}.
\li{11} Do not import a MuJoCo physics scene (Isaac Lab creates its own).
\li{12} The car is not fixed to the world (it has a free joint).
\li{13} No default density override.
\lis{14}{18} Robot type and optional actuator gain/bias overrides, all
  defaults.
\lis{19}{20} The two post-processing passes (Chapter~\ref{ch:asset}).
\li{21} Debug off.
\lis{22}{24} A trailer with the conversion date.
\end{explain}

\code{.asset_hash} holds one MD5 string computed from the converter config
and the MJCF contents. Isaac Lab compares it with a fresh hash on the next
conversion and skips the work when they match (unless
\srcpath{force_usd_conversion} is set). It is not consulted at training time.
